% Define block styles
\tikzset{
  decision/.style={
        diamond,
        draw, thick,
        text width=5.5em,
        text badly centered,
        inner sep=0pt
    },
    block/.style={
        rectangle,
        draw, thick,
        text width=10em,
        text centered,
        rounded corners
    },
    cloud/.style={
        draw,
        ellipse,
        minimum height=2em
    },
    descr/.style={
        fill=white,
        inner sep=2.5pt
    },
    connector/.style={
        -latex,
        font=\scriptsize
    },
    rectangle connector/.style={
        connector,
        to path={(\tikztostart) -- ++(#1,0pt) \tikztonodes |- (\tikztotarget) },
        pos=0.5
    },
    rectangle connector/.default=-2cm,
    straight connector/.style={
        connector,
        to path=--(\tikztotarget) \tikztonodes
    },
    line/.style={>=latex,->,thick},
    evey edge quotes/.style = {auto=right},
}

\begin{figure*}
    \centering
    \tikzsetnextfilename{algorithm}
    \begin{tikzpicture}
        \matrix (m) [matrix of nodes, column sep=1cm, row sep=6mm, align=center, nodes={anchor=center}, nodes in empty cells]    
        {
            |[block]| {Start} & &\\
            |[block]| {Generate nuclear configuration} & &\\
            |[block]| {Kick random nucleon} & &\\
            |[block]| {Evolve propagating nucleon(s) for time $\delta t$} & &\\
            |[decision]| {Possible interaction?} & |[decision]| {Outside nucleus?} & |[block]| {Remove from propagating particle list}\\
            |[block]| {Generate outgoing momenta} &  & \\
            |[decision]| {Pauli Blocked?} & |[block]|{Restore original outgoing momenta} & & \\
            |[block]| {Update outgoing momenta of particles} & &|[block]| {Add particle to list and update formation zone}& \\
        };
        \path[line] (m-1-1) edge (m-2-1);
        \path[line] (m-2-1) edge (m-3-1);
        \path[line] (m-3-1) edge (m-4-1);
        \path[line] (m-4-1) edge (m-5-1);
        \path[line] (m-5-1) edge node[midway, above, text=black]{No} (m-5-2);
        \path[line] (m-5-1) edge node[midway, right, text=black]{Yes} (m-6-1);
        \path[line] (m-5-2) edge node[midway, above, text=black]{Yes} (m-5-3);
        \draw[line] (m-5-2) |- node[midway, right, text=black]{No} (m-4-1);
        \path[line] (m-6-1) edge (m-7-1);
        \path[line] (m-7-1) edge node[midway, above, text=black]{Yes} (m-7-2);
        \path[line] (m-7-2) edge (m-5-2);
        \path[line] (m-7-1) edge node[midway, right, text=black]{No} (m-8-1);
        \path[line] (m-8-1) edge (m-8-3);
        \draw[line] (m-8-3) -- (m-7-3.center) -- (m-6-3.center) -- (m-6-2.center) -- (m-5-2);
    \end{tikzpicture}
    \caption{
        The proposed algorithm for the INC model.
        \label{fig:algorithm}
    }
\end{figure*}